\documentclass[10pt,a4paper]{amsart}
\usepackage{amsmath,amssymb,amsthm,graphicx}
\usepackage{verbatim}
\newenvironment{reference}{\comment}{\endcomment}
%\newenvironment{reference}{}{}
\newenvironment{slogan}{\comment}{\endcomment}
\newenvironment{history}{\comment}{\endcomment}

% For commutative diagrams we use Xy-pic
\usepackage[all]{xy}

% We use 2cell for 2-commutative diagrams.
\xyoption{2cell}
% Theorem environments.
%
\theoremstyle{plain}
\newtheorem{theorem}[subsection]{Theorem}
\newtheorem{proposition}[subsection]{Proposition}
\newtheorem{lemma}[subsection]{Lemma}

\theoremstyle{definition}
\newtheorem{definition}[subsection]{Definition}
\newtheorem{example}[subsection]{Example}
\newtheorem{exercise}[subsection]{Exercise}
\newtheorem{situation}[subsection]{Situation}

\theoremstyle{remark}
\newtheorem{remark}[subsection]{Remark}
\newtheorem{remarks}[subsection]{Remarks}

\numberwithin{equation}{subsection}

% Macros
%
\def\lim{\mathop{\mathrm{lim}}\nolimits}
\def\colim{\mathop{\mathrm{colim}}\nolimits}
\def\Spec{\mathop{\mathrm{Spec}}}
\def\Hom{\mathop{\mathrm{Hom}}\nolimits}
\def\Ext{\mathop{\mathrm{Ext}}\nolimits}
\def\SheafHom{\mathop{\mathcal{H}\!\mathit{om}}\nolimits}
\def\SheafExt{\mathop{\mathcal{E}\!\mathit{xt}}\nolimits}
\def\Sch{\mathit{Sch}}
\def\Mor{\mathop{\mathrm{Mor}}\nolimits}
\def\Ob{\mathop{\mathrm{Ob}}\nolimits}
\def\Sh{\mathop{\mathit{Sh}}\nolimits}
\def\NL{\mathop{N\!L}\nolimits}
\def\CH{\mathop{\mathrm{CH}}\nolimits}
\def\proetale{{pro\text{-}\acute{e}tale}}
\def\etale{{\acute{e}tale}}
\def\QCoh{\mathit{QCoh}}
\def\Ker{\mathop{\mathrm{Ker}}}
\def\Im{\mathop{\mathrm{Im}}}
\def\Coker{\mathop{\mathrm{Coker}}}
\def\Coim{\mathop{\mathrm{Coim}}}

% Boxtimes
%
\DeclareMathSymbol{\boxtimes}{\mathbin}{AMSa}{"02}

%
% Macros for moduli stacks/spaces
%
\def\QCohstack{\mathcal{QC}\!\mathit{oh}}
\def\Cohstack{\mathcal{C}\!\mathit{oh}}
\def\Spacesstack{\mathcal{S}\!\mathit{paces}}
\def\Quotfunctor{\mathrm{Quot}}
\def\Hilbfunctor{\mathrm{Hilb}}
\def\Curvesstack{\mathcal{C}\!\mathit{urves}}
\def\Polarizedstack{\mathcal{P}\!\mathit{olarized}}
\def\Complexesstack{\mathcal{C}\!\mathit{omplexes}}
% \Pic is the operator that assigns to X its picard group, usage \Pic(X)
% \Picardstack_{X/B} denotes the Picard stack of X over B
% \Picardfunctor_{X/B} denotes the Picard functor of X over B
\def\Pic{\mathop{\mathrm{Pic}}\nolimits}
\def\Picardstack{\mathcal{P}\!\mathit{ic}}
\def\Picardfunctor{\mathrm{Pic}}
\def\Deformationcategory{\mathcal{D}\!\mathit{ef}}
\usepackage{cleveref}
\crefname{originallemma}{original lemma}{original lemmas}
\Crefname{originallemma}{Original lemma}{Original lemmas}
\makeatletter
\newlabel{lemma-affine-equivalence-algebras}{{081V}{source}{}{sourcepointer-1}{}}
\newlabel{lemma-affine-equivalence-algebras@cref}{{[originallemma][0][]081V}{[1][0][]source}{}{}{}}
\makeatother
\crefname{originallemma}{original lemma}{original lemmas}
\Crefname{originallemma}{Original lemma}{Original lemmas}
\makeatletter
\newlabel{lemma-quasi-affine-local}{{03WM}{source}{}{sourcepointer-2}{}}
\newlabel{lemma-quasi-affine-local@cref}{{[originallemma][0][]03WM}{[1][0][]source}{}{}{}}
\makeatother
\crefname{originallemma}{original lemma}{original lemmas}
\Crefname{originallemma}{Original lemma}{Original lemmas}
\makeatletter
\newlabel{morphisms-lemma-characterize-quasi-affine}{{01SM}{source}{}{sourcepointer-3}{}}
\newlabel{morphisms-lemma-characterize-quasi-affine@cref}{{[originallemma][0][]01SM}{[1][0][]source}{}{}{}}
\makeatother
\crefname{originalremark}{original remark}{original remarks}
\Crefname{originalremark}{Original remark}{Original remarks}
\makeatletter
\newlabel{remark-factorization-quasi-compact-quasi-separated}{{081X}{source}{}{sourcepointer-4}{}}
\newlabel{remark-factorization-quasi-compact-quasi-separated@cref}{{[originalremark][0][]081X}{[1][0][]source}{}{}{}}
\makeatother
\crefname{originallemma}{original lemma}{original lemmas}
\Crefname{originallemma}{Original lemma}{Original lemmas}
\makeatletter
\newlabel{spaces-properties-lemma-pushforward-etale-base-change-modules}{{03LX}{source}{}{sourcepointer-5}{}}
\newlabel{spaces-properties-lemma-pushforward-etale-base-change-modules@cref}{{[originallemma][0][]03LX}{[1][0][]source}{}{}{}}
\makeatother
\title{Quasi-affineness characterized by the factorization through the relative spectrum}
\author{The Stacks Project Authors}
\date{stacks-project:spaces-morphisms.tex:local-snapshot-20261012; GNU FDL 1.2 or later}
\begin{document}
\maketitle
\setcounter{subsection}{1}
\noindent\textit{Editorial note.} All delivered mathematical text is the author's own, copied from the pinned Stacks Project chapter file spaces-morphisms.tex; no statement or proof step was rewritten, reordered or removed.\par\medskip
\noindent\textit{Editorial note.} The statement refers to the canonical factorization named in a remark of this chapter; the remark is a navigation pointer and the proof works with the factorization as displayed.\par\medskip
\noindent\textit{Editorial note.} The five references of the counted spans are recorded as non-load-bearing navigation pointers; the argument that consumes them stays inside the counted proof span.\par\medskip
\noindent\textit{Editorial note.} Preamble: the three native excerpt spans of sources/stacks-spaces-morphisms/preamble.tex (lines 8-13, 14-18 and 154-217). The chapter file itself carries no preamble text.\par\medskip
\noindent\textit{Allowed standard prerequisites (editorial).}\begin{itemize}
\item Quasi-affineness is local on the base. Lemma lemma-quasi-affine-local, used at proof lines 4281-4283. Allows the reduction to the case of a scheme base.
\item The base change formula for pushforwards of structure sheaves along an etale morphism. Properties of Spaces, Lemma spaces-properties-lemma-pushforward-etale-base-change-modules, used at proof lines 4283-4286. Identifies the restriction of f\_*O\_Y with the pushforward along the base change.
\item Formation of the relative spectrum commutes with base change. Lemma lemma-affine-equivalence-algebras, used at proof lines 4286-4289. Makes the factorization compatible with the reduction.
\item The scheme case of the characterization of quasi-affine morphisms. Morphisms, Lemma morphisms-lemma-characterize-quasi-affine, used at proof lines 4290-4293. Is the input applied once X and Y are schemes.
\end{itemize}\medskip
\ifdefined\hypertarget\hypertarget{sourcepointer-1}{}\fi
\noindent\textit{Source reference.} Original lemma 081V: focused-seeds/AGP/sources/stacks-snapshot-20261011/tags.txt, source lines 10230--10230. Navigation pointer into the Stacks Project; the standard fact it supplies is applied inside the counted proof span and the argument resting on it stays there.\par\medskip
\ifdefined\hypertarget\hypertarget{sourcepointer-2}{}\fi
\noindent\textit{Source reference.} Original lemma 03WM: focused-seeds/AGP/sources/stacks-snapshot-20261011/tags.txt, source lines 5117--5117. Navigation pointer into the Stacks Project; the standard fact it supplies is applied inside the counted proof span and the argument resting on it stays there.\par\medskip
\ifdefined\hypertarget\hypertarget{sourcepointer-3}{}\fi
\noindent\textit{Source reference.} Original lemma 01SM: focused-seeds/AGP/sources/stacks-snapshot-20261011/tags.txt, source lines 2343--2343. Navigation pointer into the Stacks Project; the standard fact it supplies is applied inside the counted proof span and the argument resting on it stays there.\par\medskip
\ifdefined\hypertarget\hypertarget{sourcepointer-4}{}\fi
\noindent\textit{Source reference.} Original remark 081X: focused-seeds/AGP/sources/stacks-snapshot-20261011/tags.txt, source lines 10232--10232. Navigation pointer into the Stacks Project; the standard fact it supplies is applied inside the counted proof span and the argument resting on it stays there.\par\medskip
\ifdefined\hypertarget\hypertarget{sourcepointer-5}{}\fi
\noindent\textit{Source reference.} Original lemma 03LX: focused-seeds/AGP/sources/stacks-snapshot-20261011/tags.txt, source lines 4729--4729. Navigation pointer into the Stacks Project; the standard fact it supplies is applied inside the counted proof span and the argument resting on it stays there.\par\medskip
\begin{lemma}
\label{lemma-characterize-quasi-affine}
Let $S$ be a scheme.
A quasi-compact and quasi-separated morphism of algebraic spaces
$f : Y \to X$ is quasi-affine if and only if the canonical factorization
$Y \to \underline{\Spec}_X(f_*\mathcal{O}_Y)$
(Remark \ref{remark-factorization-quasi-compact-quasi-separated})
is an open immersion.
\end{lemma}

\begin{proof}
Let $U \to X$ be a surjective morphism where $U$ is a scheme.
Since we may check whether $f$ is quasi-affine after base change to
$U$ (Lemma \ref{lemma-quasi-affine-local}), since $f_*\mathcal{O}_Y|_V$
is equal to $(Y \times_X U \to U)_*\mathcal{O}_{Y \times_X U}$
(Properties of Spaces, Lemma
\ref{spaces-properties-lemma-pushforward-etale-base-change-modules}), and
since formation of relative spectrum commutes with base change
(Lemma \ref{lemma-affine-equivalence-algebras}),
we see that the assertion reduces to the case that $X$ is a scheme.
If $X$ is a scheme and either $f$ is quasi-affine or
$Y \to \underline{\Spec}_X(f_*\mathcal{O}_Y)$ is an open immersion,
then $Y$ is a scheme as well. Thus we reduce to
Morphisms, Lemma \ref{morphisms-lemma-characterize-quasi-affine}.
\end{proof}

\end{document}
